(1) 在各向同性低场 Drude 模型下，每类载流子的纵向电导分别为 \(ne\mu_e\) 与 \(pe\mu_h\)，均为正数，这是因为电子漂移速度与电场反向但电荷也为负。零场 \(j_x=e(n\mu_e+p\mu_h)E_x\)，故 \(\sigma_0=e(n\mu_e+p\mu_h)\)。忽略两载流子间相互拖曳、迁移率场依赖和高阶磁场项。
(2) 将每类载流子的稳态洛伦兹力平衡求到 \(B\) 的一次项，电子和空穴的横向电导贡献符号相反。施加开路条件 \(j_y=0\)，消去横向电场，得到 \(R_H=(p\mu_h^2-n\mu_e^2)/[e(n\mu_e+p\mu_h)^2]\)。若只有电子，\(p=0\)，化为 \(-1/(ne)\)；若只有空穴，\(n=0\)，化为 \(+1/(pe)\)，与规定的符号自洽。
(3) 霍尔电压符号在 \(p\mu_h^2=n\mu_e^2\) 时翻转，而非简单在 \(p=n\) 时翻转。较少数但迁移率更高的载流子可以在横向响应中占优势，因为洛伦兹横向漂移相对于一次电场响应再多出一个迁移率因子。这是以霍尔测量推断多数载流子类型时需要辨别的双载流子效应。
(4) 由 Drude 张量分别展开 \(\sigma_{xx}\)、\(\sigma_{xy}\) 并逆矩阵，固定 \(j_y=0\) 得 \(\rho_{xx}(B)/\rho_{xx}(0)=1+\{np\mu_e\mu_h(\mu_e+\mu_h)^2/(n\mu_e+p\mu_h)^2\}B^2+O(B^4)\)。系数非负；当 \(n=0\) 或 \(p=0\) 时二阶磁阻消失。这个结果只适用于各载流子为单一弛豫时间且 \(\mu_{e,h}|B|\ll1\) 的理想模型。
